\begin{tikzpicture}[font=\small,line cap=round,line join=round]\draw[otNavy,line width=.8pt](0,0)circle(1.7);\fill[otNavy](0,0)circle(1.4pt)node[below,font=\small,text=otNavy]{$I$};\draw[geoTrait](-0.718,1.541)--(-1.700,0.000)--(1.700,0.000)--(-0.718,1.541);\draw[otMuted,line width=.6pt](-0.810,1.397)--(-0.666,1.306)--(-0.575,1.449);\draw[otNavy,line width=.7pt](-0.850,-0.110)--(-0.850,0.110);\draw[otNavy,line width=.7pt](0.850,-0.110)--(0.850,0.110);\fill[otNavy](-0.718,1.541)circle(1.3pt)node[above,geoLabel,font=\small]{$ A $};\fill[otNavy](-1.700,0.000)circle(1.3pt)node[left,geoLabel,font=\small]{$ B $};\fill[otNavy](1.700,0.000)circle(1.3pt)node[right,geoLabel,font=\small]{$ C $};\draw[otNavy,line width=.6pt](-0.718,1.541)--(0.000,0.000);\draw[otNavy,line width=.7pt](-0.459,0.724)--(-0.260,0.817);\end{tikzpicture}\par\vspace{1mm}{\fontsize{8}{10}\selectfont\color{otMuted}Le cercle circonscrit au triangle ABC, rectangle en A, a pour diamètre l’hypoténuse $[BC]$.}